
\documentclass{article}
\usepackage{colortbl}
\usepackage{makecell}
\usepackage{multirow}
\usepackage{supertabular}

\begin{document}

\newcounter{utterance}

\centering \large Interaction Transcript for game `eqbench', experiment `eq\_bench', episode 71 with reflexion{-}Qwen3.5{-}27B.
\vspace{24pt}

{ \footnotesize  \setcounter{utterance}{1}
\setlength{\tabcolsep}{0pt}
\begin{supertabular}{c@{$\;$}|p{.15\linewidth}@{}p{.15\linewidth}p{.15\linewidth}p{.15\linewidth}p{.15\linewidth}p{.15\linewidth}}
    \# & \multicolumn{2}{c}{Player} && \multicolumn{2}{c}{Game Master} \\
    \hline

    \theutterance \stepcounter{utterance}  
    & & & \multicolumn{4}{p{0.6\linewidth}}{
        \cellcolor[rgb]{0.9,0.9,0.9}{
            \makecell[{{p{\linewidth}}}]{
                \texttt{\tiny{[P1$\langle$GM]}}
                \texttt{You have attempted this same task before and did not succeed. These are the plans you wrote after those attempts.} \\
\texttt{Trial \#1: New plan: My previous attempt failed because I assigned a non{-}zero score to "Anxious" (3) and "Impatient" (2), which contradicts Richard's character arc in the dialogue. Richard is portrayed as manipulative, self{-}assured, and in control, using phrases like "I trust your judgment" and "strategic communication" to mask his agenda. He shows no signs of nervousness or frustration. Next time, I will strictly align the scores with his dominant trait of arrogance and control: I will assign a high score (e.g., 9 or 10) to "Confident" and set "Anxious," "Impatient," and "Ashamed" to 0, as there is no textual evidence supporting those emotions.} \\
\\ 
\texttt{Trial \#2: New plan: My previous attempt failed because I assigned a non{-}zero score to "Anxious" (3) and "Impatient" (2), which contradicts Richard's character arc in the dialogue. Richard is portrayed as manipulative, self{-}assured, and in control, using phrases like "I trust your judgment" and "strategic communication" to mask his agenda. He shows no signs of nervousness or frustration. Next time, I will strictly align the scores with his dominant trait of arrogance and control: I will assign a high score (e.g., 9 or 10) to "Confident" and set "Anxious," "Impatient," and "Ashamed" to 0, as there is no textual evidence supporting those emotions.} \\
\\ 
\\ 
\\ 
\texttt{Your task is to predict the likely emotional responses of a character in this dialogue:} \\
\\ 
\texttt{Richard: Emily, always a pleasure. You're looking... tenacious as always.} \\
\texttt{Emily: Cut the flattery, Richard. What do you want?} \\
\texttt{Richard: Straight to the point. I like that. I have a story for you. A big one.} \\
\texttt{Emily: I'm listening.} \\
\texttt{Richard: It's about my opponent. There are some... unsavory details about his past that the public deserves to know.} \\
\texttt{Emily: And you're sharing this out of the goodness of your heart?} \\
\texttt{Richard: Let's just say, I believe in transparency.} \\
\texttt{Emily: You mean manipulation.} \\
\texttt{Richard: That's a harsh word, Emily. I prefer 'strategic communication'.} \\
\texttt{Emily: Call it what you want, Richard. I'm not your puppet.} \\
\texttt{Richard: Of course not. I respect your work, Emily. That's why I came to you.} \\
\texttt{Emily: I'm sure.} \\
\texttt{Richard: I have evidence. It's in this envelope.} \\
\texttt{Emily: And you expect me to just publish this without verifying?} \\
\texttt{Richard: I trust your judgment, Emily. You'll see it's all true.} \\
\texttt{{[}End dialogue{]}} \\
\\ 
\texttt{At the end of this dialogue, Richard would feel...} \\
\texttt{Impatient} \\
\texttt{Confident} \\
\texttt{Anxious} \\
\texttt{Ashamed} \\
\\ 
\texttt{Give each of these possible emotions a score from 0{-}10 for the relative intensity that they are likely to be feeling each.} \\
\\ 
\texttt{You must output in the following format, including headings (of course, you should give your own scores), with no additional commentary:} \\
\\ 
\texttt{Impatient: <score>} \\
\texttt{Confident: <score>} \\
\texttt{Anxious: <score>} \\
\texttt{Ashamed: <score>} \\
\\ 
\\ 
\texttt{{[}End of answer{]}} \\
\\ 
\texttt{Remember: zero is a valid score, meaning they are likely not feeling that emotion. You must score at least one emotion > 0.} \\
\\ 
\texttt{Your answer:} \\
            }
        }
    }
     \\ \\

    \theutterance \stepcounter{utterance}  
    & \multicolumn{4}{p{0.6\linewidth}}{
        \cellcolor[rgb]{0.561,0.737,0.561}{
            \makecell[{{p{\linewidth}}}]{
                \texttt{\tiny{[P1$\rangle$GM]}}
                \texttt{Impatient: 0} \\
\texttt{Confident: 10} \\
\texttt{Anxious: 0} \\
\texttt{Ashamed: 0} \\
            }
        }
    }
    & & \\ \\

    \theutterance \stepcounter{utterance}  
    & & & \multicolumn{2}{p{0.3\linewidth}}{
        \cellcolor[rgb]{0.9,0.9,0.9}{
            \makecell[{{p{\linewidth}}}]{
                \texttt{\tiny{[GM$|$GM]}}
                \texttt{Impatient: 0} \\
\texttt{Confident: 10} \\
\texttt{Anxious: 0} \\
\texttt{Ashamed: 0} \\
            }
        }
    }
    & & \\ \\

    \theutterance \stepcounter{utterance}  
    & & & \multicolumn{2}{p{0.3\linewidth}}{
        \cellcolor[rgb]{0.9,0.9,0.9}{
            \makecell[{{p{\linewidth}}}]{
                \texttt{\tiny{[GM$|$GM]}}
                \texttt{Impatient: 3} \\
\texttt{Confident: 8} \\
\texttt{Anxious: 2} \\
\texttt{Ashamed: 0} \\
            }
        }
    }
    & & \\ \\

    \theutterance \stepcounter{utterance}  
    & & & \multicolumn{2}{p{0.3\linewidth}}{
        \cellcolor[rgb]{0.9,0.9,0.9}{
            \makecell[{{p{\linewidth}}}]{
                \texttt{\tiny{[GM$|$GM]}}
                \texttt{game\_result = LOSE} \\
            }
        }
    }
    & & \\ \\

    \theutterance \stepcounter{utterance}  
    & & \multicolumn{4}{p{0.6\linewidth}}{
        \cellcolor[rgb]{0.95,0.95,0.95}{
            \makecell[{{p{\linewidth}}}]{
                \texttt{\tiny{[P1$\leftrightarrow$P]}}
                \texttt{Post{-}episode Reflexion (not sent to GM):} \\
\texttt{New plan: My previous attempt failed because I assigned a score of 0 to "Impatient," overlooking the textual evidence of Richard's frustration. Although he is confident, his repeated attempts to persuade Emily ("I trust your judgment," "You'll see it's all true") after her sharp rejections ("Cut the flattery," "I'm not your puppet") indicate a growing desire to move the interaction forward, which aligns with impatience. Next time, I will assign a high score to "Confident" (e.g., 9) to reflect his self{-}assured manipulation, but I will also assign a moderate score to "Impatient" (e.g., 4 or 5) to account for his eagerness to bypass her skepticism. I will keep "Anxious" and "Ashamed" at 0, as there is no evidence of fear or guilt in his dialogue.} \\
            }
        }
    }
    & \\ \\

\end{supertabular}
}

\end{document}
